\documentclass[12pt,a4paper]{article}

\usepackage[margin=1in]{geometry}
\usepackage{amsmath}
\usepackage{amssymb}
\usepackage{graphicx}
\usepackage{booktabs}
\usepackage{hyperref}
\usepackage{array}
\usepackage{multirow}
\usepackage{float}
\usepackage{xcolor}
\usepackage{enumitem}
\usepackage{titlesec}
\usepackage{caption}
\usepackage[numbers]{natbib}

\hypersetup{
    colorlinks=true,
    linkcolor=blue,
    citecolor=blue,
    urlcolor=blue
}

\titleformat{\section}{\large\bfseries}{\thesection}{1em}{}
\titleformat{\subsection}{\normalsize\bfseries}{\thesubsection}{1em}{}

\title{
    \vspace{-1cm}
    \textbf{Visual Product Search Engine} \\
    \large Visual Recognition Course -- Final Project Report
}

\author{
    Arnav Oruganty (IMT2023078) \quad
    Ankith Kini (IMT2023075) \\
    Dhruv Joshi (IMT2023032) \quad
    Mithilesh Sai Yechuri (IMT2023507)
}

\date{\today}

\begin{document}

\maketitle
\tableofcontents
\newpage

% ─────────────────────────────────────────────────────────────────────────────
\section{Problem Motivation}
% ─────────────────────────────────────────────────────────────────────────────

Online shopping platforms primarily rely on keyword-based search systems.
However, users often encounter situations where they see a clothing item
they like but cannot accurately describe it using text alone.

This creates two major problems:

\begin{itemize}[noitemsep]
    \item \textbf{Language mismatch:} Users may search using simple keywords
          such as ``black hoodie'', while catalog descriptions use technical
          or brand-specific terminology.
    \item \textbf{Inconsistent metadata:} Similar products may be described
          differently across sellers, making text-based retrieval unreliable.
\end{itemize}

To solve this problem, we build a \textbf{query-by-image fashion retrieval
system}. Instead of relying only on text search, the user uploads an image,
and the system retrieves visually and semantically similar products from a
large fashion catalog.

We use the \textbf{DeepFashion In-Shop Clothes Retrieval} dataset, where
images sharing the same \texttt{item\_id} are treated as the same product
during training and evaluation.

The proposed system contains two main stages:

\begin{itemize}[noitemsep]
    \item \textbf{Offline indexing pipeline:} Generates and stores embeddings
          for all gallery products.
    \item \textbf{Online retrieval pipeline:} Processes a user query image
          and retrieves the most similar products efficiently.
\end{itemize}

% ─────────────────────────────────────────────────────────────────────────────
\section{Dataset}
% ─────────────────────────────────────────────────────────────────────────────

\subsection{DeepFashion In-Shop Clothes Retrieval}

The DeepFashion In-Shop Clothes Retrieval benchmark contains high-resolution
fashion images grouped according to item identity. Each product contains
multiple images captured from different viewpoints and poses.

\begin{table}[H]
\centering
\caption{Dataset split statistics.}
\begin{tabular}{lrr}
\toprule
\textbf{Split} & \textbf{Images} & \textbf{Role} \\
\midrule
Train   & 25,882 & Fine-tune CLIP \\
Query   & 14,218 & Query images for retrieval \\
Gallery & 12,612 & Search database \\
\midrule
Total   & 52,712 & \\
\bottomrule
\end{tabular}
\end{table}

The dataset additionally provides the annotation file \texttt{list\_bbox\_inshop.txt}, which contains ground-truth bounding box coordinates for all images. These annotations are used during training and evaluation to crop clothing regions and compute retrieval and captioning metrics under controlled settings.

However, during inference, ground-truth bounding boxes are unavailable for unseen images. Therefore, the bounding boxes are predicted by a YOLO model.

\subsection{Ground Truth}

Two images are considered a correct retrieval match if and only if they
share the same \texttt{item\_id}. Since each clothing item in the
DeepFashion In-Shop dataset is captured under multiple viewpoints,
poses, scales, and lighting conditions, the shared \texttt{item\_id}
serves as the identity label for evaluating retrieval performance.
During evaluation, all gallery images belonging to the same
\texttt{item\_id} as the query image are treated as relevant matches,
while images associated with different item identities are considered
non-relevant.

\subsection{Preprocessing}

For training and evaluation, we use the ground-truth bounding box
annotations provided in \texttt{list\_bbox\_inshop.txt} to crop the
primary clothing regions from each image. Using the annotated bounding
boxes reduces background clutter and ensures that the retrieval model
focuses primarily on the garment itself.

We also experimented with using YOLOv8-generated detections for both
training and evaluation in order to simulate a fully automated pipeline.
However, we observed that retrieval performance using ground-truth crops
was consistently better, likely because the annotations provide tighter
and more accurate localization of the clothing items.

For the final inference pipeline on unseen query images, where
ground-truth annotations are unavailable, we use a YOLOv8 detector to
automatically localize and crop the primary clothing region. If YOLO
fails to detect any object, the original full image is used as a
fallback.

The cropped images and associated metadata are stored offline for
efficient retrieval.

% ─────────────────────────────────────────────────────────────────────────────
\section{System Architecture}
% ─────────────────────────────────────────────────────────────────────────────

The proposed system consists of two components:

\begin{itemize}[noitemsep]
    \item Offline indexing pipeline: prepare product embeddings and store them in a vector index
    \item Online retrieval pipeline: process a query image and return top matching products quickly
\end{itemize}

\subsection{Offline Indexing Pipeline}

For each gallery image, the following operations are performed:

\begin{enumerate}[noitemsep]
    \item \textbf{Product localization (YOLO).}

    During training and evaluation, we use the ground-truth bounding box
    annotations from \texttt{list\_bbox\_inshop.txt} to crop the primary
    clothing region.
    For the final deployment pipeline and Streamlit demo, 
    YOLOv8 is used to automatically detect and crop the clothing item.

    \item \textbf{Semantic captioning (BLIP-2).}

    BLIP-2 generates a semantic caption describing clothing attributes such
    as color, fit, texture, and style.

    \item \textbf{Cross-modal embedding generation (CLIP).}

    The cropped image and generated caption are independently encoded using
    CLIP's visual and text encoders.

    The final embedding is obtained using weighted fusion:

    \begin{equation}
        \mathbf{v}_i =
        \alpha \phi_V(\hat{x}_i)
        +
        (1-\alpha)\phi_T(c_i),
        \qquad
        \|\mathbf{v}_i\| = 1
    \end{equation}

    where:

    \begin{itemize}[noitemsep]
        \item $\hat{x}_i$ is the cropped image,
        \item $c_i$ is the generated caption,
        \item $\phi_V(\cdot)$ is the CLIP visual encoder,
        \item $\phi_T(\cdot)$ is the CLIP text encoder,
        \item $\alpha \in [0,1]$ controls image--text fusion.
    \end{itemize}

    \item \textbf{ANN indexing using HNSW.}

    The final embeddings are stored inside a Hierarchical Navigable Small World (HNSW) index using \texttt{hnswlib} for fast approximate nearest-neighbor retrieval.
\end{enumerate}

\subsection{Online Query Pipeline}

Given a query image $q$, retrieval proceeds as follows:

\begin{enumerate}[noitemsep]
    \item \textbf{Query preprocessing.}

    During training and evaluation, we use the ground-truth bounding box annotations
from \texttt{list\_bbox\_inshop.txt} to crop the primary clothing region. For the final
deployment pipeline and Streamlit demo, YOLOv8 is used to automatically detect
and crop the clothing item.

    \item \textbf{Query encoding.}

    The cropped query image is encoded using the CLIP visual encoder.

    \item \textbf{Candidate retrieval.}

    The HNSW index retrieves the top-$K$ nearest neighbors using cosine
    similarity.

    \item \textbf{Semantic re-ranking.}

    BLIP-2 image-text matching scores are computed between the query image
    and candidate captions to improve ranking quality.
\end{enumerate}

% ─────────────────────────────────────────────────────────────────────────────
\section{Fine-Tuning Strategy}
% ─────────────────────────────────────────────────────────────────────────────

We use pretrained YOLOv8, BLIP-2, and CLIP models.

Only CLIP is fine-tuned in this project. YOLO and BLIP-2 remain frozen.

\subsection{CLIP Fine-tuning}

\textbf{Goal.}

The objective is to pull embeddings belonging to the same
\texttt{item\_id} closer together while pushing embeddings of different
products farther apart.

\textbf{Updated parameters.}

\begin{itemize}[noitemsep]
    \item Last 4 transformer blocks of the CLIP vision encoder
    \item Projection head
\end{itemize}

The CLIP text encoder remains frozen.

\textbf{Loss function.}

We use symmetric InfoNCE contrastive loss.

Given anchor-positive pairs $(a_i,p_i)$ sampled from the same
\texttt{item\_id}:

\begin{equation}
\mathcal{L}
=
-\frac{1}{2N}
\sum_{i=1}^{N}
\left[
\log
\frac{
e^{\mathbf{a}_i \cdot \mathbf{p}_i / \tau}
}{
\sum_j e^{\mathbf{a}_i \cdot \mathbf{p}_j / \tau}
}
+
\log
\frac{
e^{\mathbf{p}_i \cdot \mathbf{a}_i / \tau}
}{
\sum_j e^{\mathbf{p}_i \cdot \mathbf{a}_j / \tau}
}
\right]
\end{equation}

where $\tau = 0.07$ is the temperature parameter.

All embeddings are L2-normalized.

\subsection{Training Details}

\begin{itemize}[noitemsep]
    \item Optimizer: AdamW
    \item Learning rate: $1 \times 10^{-5}$
    \item Batch size: 64
    \item Epochs: 5
    \item Mixed precision training using FP16
\end{itemize}

% ─────────────────────────────────────────────────────────────────────────────
\section{Evaluation Protocol}
% ─────────────────────────────────────────────────────────────────────────────

\subsection{Metrics}

We evaluate the retrieval system using standard information retrieval metrics.

\begin{itemize}[noitemsep]
    \item \textbf{Recall@K}

    Fraction of queries for which at least one correct product appears
    within the top-$K$ retrieved results.
    \begin{equation}
    \text{Recall@}K
    = \frac{\text{number of relevant items in top-}K}
           {\text{total number of relevant items}}
\end{equation}

    \item \textbf{NDCG@K}

    Position-aware ranking metric that rewards relevant products appearing
    earlier in the ranked list.
    \begin{equation}
    \text{DCG@}K = \sum_{i=1}^{K}
        \frac{\mathrm{rel}(i)}{\log_{2}(i+1)},
    \qquad
    \text{NDCG@}K = \frac{\text{DCG@}K}{\text{IDCG@}K}
\end{equation}
where $\mathrm{rel}(i) \in \{0,1\}$ is the binary relevance of the item at
rank $i$, and IDCG@$K$ is the DCG of the ideal ranking.
NDCG@$K \in [0,1]$, with 1 indicating a perfect ranking.


    \item \textbf{mAP@K}

    Mean Average Precision up to $K$, measuring both retrieval accuracy
    and ranking quality.
    For a single query, Average Precision (AP) at $K$ is:
\begin{equation}
    \text{AP@}K
    = \frac{1}{N}\sum_{k=1}^{K} \text{Precision}(k)\cdot\mathrm{rel}(k)
\end{equation}
where $N$ is the total number of relevant items, $\text{Precision}(k)$ is
the precision at rank $k$, and $\mathrm{rel}(k)=1$ if the item at position
$k$ is relevant. MAP@$K$ is then the mean of AP@$K$ over all $Q$ queries:
\begin{equation}
    \text{mAP@}K = \frac{1}{Q}\sum_{q=1}^{Q}\text{AP}_q\text{@}K
\end{equation}
MAP@$K \in [0,1]$, with 1 corresponding to an ideal ranking.
\end{itemize}

We report metrics at:

\[
K \in \{5,10,15\}
\]

\subsection{Reporting}

All results are reported as mean $\pm$ standard deviation.
Since configurations A and B involve no trainable parameters,
they are deterministic and require no seeds.
For configuration C, which involves fine-tuning CLIP,
results are averaged over 3 runs using team members'
roll numbers as random seeds; the seeding strategy
is discussed in Section 6.4.3.
% ─────────────────────────────────────────────────────────────────────────────
\section{Experiments: Ablation Study}
% ─────────────────────────────────────────────────────────────────────────────

\subsection{Experimental Conditions}

We evaluate three conditions to isolate the contribution of each component.

\begin{table}[H]
\centering
\caption{Ablation study configurations.}
\begin{tabular}{clp{7cm}}
\toprule
\textbf{ID} & \textbf{Configuration} & \textbf{Purpose} \\
\midrule
A &
Vision-only CLIP ($\alpha = 1$)
&
Baseline without semantic fusion or fine-tuning
\\

B &
Frozen CLIP + frozen BLIP-2
&
Measures gain from captioning without
fine-tuning
\\

C &
Fine-tuned CLIP + frozen BLIP-2
&
Evaluate effect of contrastive fine-tuning
\\

\bottomrule
\end{tabular}
\end{table}

All conditions use the same HNSW index configuration to ensure fair
comparison.

For Conditions B and C, we evaluate two fusion weights:

\[
\alpha = 0.7
\qquad
\text{and}
\qquad
\alpha = 0.5
\]

% ─────────────────────────────────────────────────────────────────────────────
\subsection{Results}
% ─────────────────────────────────────────────────────────────────────────────

\begin{table}[H]
\centering
\caption{Recall@K results (mean $\pm$ std, 3 seeds). Condition C entries are averaged over seeds 32, 75, and 78.}
\label{tab:recall}
\resizebox{\textwidth}{!}{
\begin{tabular}{lccc}
\toprule
\textbf{Condition} & \textbf{Recall@5} & \textbf{Recall@10} & \textbf{Recall@15} \\
\midrule
A: Vision-only CLIP ($\alpha=1$, pre-trained)
& $0.4855 \pm 0.0000$
& $0.5643 \pm 0.0000$
& $0.6093 \pm 0.0000$ \\
B: Frozen CLIP + text ($\alpha=0.5$)
& $0.3979 \pm 0.0000$ & $0.4721 \pm 0.0000$ & $0.5123 \pm 0.0000$ \\
B: Frozen CLIP + text ($\alpha=0.7$)
& $0.4247 \pm 0.0000$ & $0.4935 \pm 0.0000$ & $0.5347 \pm 0.0000$ \\
C: Fine-tuned CLIP, vision-only ($\alpha=1$)
& $0.5496 \pm 0.0009$
& $0.6189 \pm 0.0014$
& $0.6578 \pm 0.0030$ \\
C: Fine-tuned CLIP + text ($\alpha=0.5$)
& $0.5310 \pm 0.0088$
& $0.6046 \pm 0.0090$
& $0.6488 \pm 0.0083$ \\
C: Fine-tuned CLIP + text ($\alpha=0.7$)
& $\mathbf{0.5828 \pm 0.0009}$
& $\mathbf{0.6514 \pm 0.0011}$
& $\mathbf{0.6878 \pm 0.0016}$ \\
\bottomrule
\end{tabular}}
\end{table}

\begin{table}[H]
\centering
\caption{NDCG@K results (mean $\pm$ std, 3 seeds).}
\label{tab:ndcg}
\resizebox{\textwidth}{!}{
\begin{tabular}{lccc}
\toprule
\textbf{Condition} & \textbf{NDCG@5} & \textbf{NDCG@10} & \textbf{NDCG@15} \\
\midrule
A: Vision-only CLIP ($\alpha=1$, pre-trained)
& $0.5151 \pm 0.0000$
& $0.5963 \pm 0.0000$
& $0.6428 \pm 0.0000$ \\
B: Frozen CLIP + text ($\alpha=0.5$)
& $0.3605 \pm 0.0000$ & $0.4034 \pm 0.0000$ & $0.4198 \pm 0.0000$ \\
B: Frozen CLIP + text ($\alpha=0.7$)
& $0.4055 \pm 0.0000$ & $0.4475 \pm 0.0000$ & $0.4694 \pm 0.0000$ \\
C: Fine-tuned CLIP, vision-only ($\alpha=1$)
& $0.6185 \pm 0.0023$
& $0.6988 \pm 0.0033$
& $0.7446 \pm 0.0041$ \\
C: Fine-tuned CLIP + text ($\alpha=0.5$)
& $0.5856 \pm 0.0108$
& $0.6591 \pm 0.0133$
& $0.7006 \pm 0.0140$ \\
C: Fine-tuned CLIP + text ($\alpha=0.7$)
& $\mathbf{0.6603 \pm 0.0017}$
& $\mathbf{0.7396 \pm 0.0031}$
& $\mathbf{0.7853 \pm 0.0034}$ \\
\bottomrule
\end{tabular}}
\end{table}

\begin{table}[H]
\centering
\caption{mAP@K results (mean $\pm$ std, 3 seeds).}
\label{tab:map}
\resizebox{\textwidth}{!}{
\begin{tabular}{lccc}
\toprule
\textbf{Condition} & \textbf{mAP@5} & \textbf{mAP@10} & \textbf{mAP@15} \\
\midrule
A: Vision-only CLIP ($\alpha=1$, pre-trained)
& $0.3582 \pm 0.0000$
& $0.3498 \pm 0.0000$
& $0.3394 \pm 0.0000$ \\
B: Frozen CLIP + text ($\alpha=0.5$)
& $0.2442 \pm 0.0000$ & $0.2276 \pm 0.0000$ & $0.2106 \pm 0.0000$ \\
B: Frozen CLIP + text ($\alpha=0.7$)
& $0.2777 \pm 0.0000$ & $0.2571 \pm 0.0000$ & $0.2422 \pm 0.0000$ \\
C: Fine-tuned CLIP, vision-only ($\alpha=1$)
& $0.6741 \pm 0.0037$
& $0.7563 \pm 0.0052$
& $0.7969 \pm 0.0062$ \\
C: Fine-tuned CLIP + text ($\alpha=0.5$)
& $0.6330 \pm 0.0128$
& $0.7025 \pm 0.0160$
& $0.7364 \pm 0.0167$ \\
C: Fine-tuned CLIP + text ($\alpha=0.7$)
& $\mathbf{0.7211 \pm 0.0029}$
& $\mathbf{0.8011 \pm 0.0050}$
& $\mathbf{0.8411 \pm 0.0057}$ \\
\bottomrule
\end{tabular}}
\end{table}

% ─────────────────────────────────────────────────────────────────────────────

\subsection{Metric Trends}

Across all conditions, Recall@$K$ increases monotonically with $K$ because a
larger candidate pool raises the probability of including at least one correct
match. In Condition A, Recall rises from $0.4855$ at $K{=}5$ to $0.6093$ at
$K{=}15$; Condition B ($\alpha{=}0.7$) follows a similar pattern from $0.4247$
to $0.5347$; and Condition C ($\alpha{=}0.7$) achieves the strongest coverage,
reaching $0.6878$ at $K{=}15$.

NDCG@$K$ also increases with $K$, as additional relevant items beyond rank 1
contribute discounted gain. However, the marginal improvement diminishes as $K$
grows, since items at lower ranks receive smaller logarithmic weights. Condition
A improves from $0.5151$ to $0.6428$; Condition B ($\alpha{=}0.7$) from
$0.4055$ to $0.4694$; and Condition C ($\alpha{=}0.7$) from $0.6603$ to
$0.7853$.

mAP@$K$ behaves differently from Recall and NDCG: it decreases as $K$
increases across Conditions A and B. This occurs because mAP averages precision
over all relevant ranks up to $K$; introducing more candidates at lower
positions without corresponding true matches reduces the average precision.
In Condition A, mAP falls from $0.3582$ at $K{=}5$ to $0.3394$ at $K{=}15$,
and in Condition B ($\alpha{=}0.7$) from $0.2777$ to $0.2422$. Condition C
reverses this trend entirely - mAP rises from $0.7211$ to $0.8411$ at
$\alpha{=}0.7$ - because the fine-tuned encoder places correct matches
consistently at rank 1, so higher $K$ includes additional true positives that
sustain precision rather than diluting it.

% ─────────────────────────────────────────────────────────────────────────────
\subsection{Detailed Case-wise Analysis}
% ─────────────────────────────────────────────────────────────────────────────

\subsubsection{Condition A: Vision-only CLIP ($\alpha = 1$)}

Condition A serves as the baseline configuration where only the CLIP visual
encoder $\phi_V$ is used for retrieval, with $\alpha = 1$. Formally, the
fused embedding reduces to a pure visual embedding:
\begin{equation}
    \mathbf{v}_i = \phi_V(\hat{x}_i), \qquad \|\mathbf{v}_i\| = 1
\end{equation}
No semantic caption information is incorporated and the CLIP encoder is used
without any task-specific fine-tuning. The retrieval process therefore depends
entirely on the pretrained CLIP visual embedding space.

Although CLIP provides strong general visual representations, it is not
optimized for fine-grained fashion retrieval tasks. The pretrained embedding
space tends to cluster images by broad visual properties such as overall color
and silhouette, rather than by product identity. As a result, visually similar
products with different \texttt{item\_id} values may still appear close
together in the embedding space. Typical retrieval errors include confusion
between:
\begin{itemize}[noitemsep]
    \item hoodies and sweatshirts,
    \item formal shirts and casual shirts,
    \item different sleeve styles,
    \item and similar-colored garments of different types.
\end{itemize}

The Recall@$K$ scores are therefore expected to be the lowest among all
conditions, since the model has no mechanism to leverage semantic or
identity-specific information. Similarly, NDCG@$K$ and mAP@$K$ remain
relatively weak because relevant products often appear lower in the ranked
list rather than in the top positions. Overall, Condition A establishes the
lower-bound baseline retrieval capability of generic CLIP embeddings without
semantic alignment or task-specific learning, against which the gains from
text fusion (Condition B) and contrastive fine-tuning (Condition C) are
measured.
\vspace{0.3cm}

\subsubsection{Condition B: Frozen CLIP + BLIP-2 Text Fusion}
\label{sec:condition_b}

Condition B introduces BLIP-2-generated semantic captions into the retrieval
pipeline without any task-specific fine-tuning. Gallery embeddings are formed
by fusing the frozen CLIP visual embedding $\phi_V(\hat{x}_i)$ with the
frozen CLIP text embedding of the BLIP-2 caption $\phi_T(c_i)$ at two fusion
weights, $\alpha \in \{0.5, 0.7\}$. The query embedding remains purely visual
($\alpha_{\text{query}} = 1$), as only gallery vectors incorporate caption
information.

\paragraph{Implementation.}
Captions were precomputed using BLIP-2, with image-text matching scores
computed via \texttt{BlipForImageTextRetrieval} (\texttt{Salesforce/blip2-itm-vit-g})
which provides the ITM head used for re-ranking. Gallery visual and caption embeddings of dimension $768$ were loaded
from precomputed \texttt{.npy} files (shapes $12612 \times 768$ each), verified
for filename alignment, fused according to the equation, and
L2-normalized before insertion into the \texttt{hnswlib} HNSW index.
Each query retrieved the top-15
candidates by cosine similarity, which were then re-ranked using BLIP-2 ITM
scores computed in a single batched pass, with separate re-ranked lists stored
for $K \in \{5, 10, 15\}$.

\paragraph{Performance relative to Condition A.}
Introducing text fusion without fine-tuning \emph{degrades} retrieval
performance relative to the vision-only Condition A baseline across all metrics
and both $\alpha$ values. At $\alpha{=}0.5$, Recall@15 falls from $0.6093$
(A) to $0.5123$; NDCG@5 drops from $0.5151$ to $0.3605$; and mAP@5 drops from
$0.3582$ to $0.2442$. The degradation is less severe at $\alpha{=}0.7$, where
Recall@15 recovers partially to $0.5347$ and NDCG@5 to $0.4055$, but remains
below Condition A on every metric.

\paragraph{Why caption fusion hurts without fine-tuning.}
The CLIP visual encoder produces embeddings that capture broad visual similarity
in a general-purpose space. The BLIP-2 text encoder, also frozen, produces
caption embeddings that are \emph{category-level}: the text encoding of
``a white linen shirt with button collar'' is nearly identical for every white
shirt in the gallery, regardless of which specific product it belongs to.
Fusing this category-level signal into the visual embedding collapses the
fine-grained geometric distinctions that allow the visual encoder to
discriminate between similar products. At $\alpha{=}0.5$, the two modalities
contribute equally, and the caption signal overwhelms the instance-specific
visual geometry, causing the fused vectors of visually distinct but
semantically similar items to cluster together. At $\alpha{=}0.7$, the visual
component remains dominant, limiting the damage — but the caption embedding
still injects enough semantic noise to reduce separability within the same
broad category. Since neither encoder has been aligned to the product-identity
retrieval task, the fusion produces a representational mismatch rather than a
complementary signal.


Condition B establishes that semantic captions alone, without task-specific
alignment, are insufficient to improve retrieval and can actively harm
performance by diluting instance-discriminative visual signals. The gain from
captioning only materialises when the visual encoder is fine-tuned to produce
identity-discriminative embeddings — as demonstrated in Condition C.

\subsubsection{Condition C: Fine-tuned CLIP + BLIP-2 Text Fusion}

Condition C adds task-specific contrastive fine-tuning on top of the
multimodal fusion framework established in Condition B.
The hypothesis is that a vision encoder explicitly trained to pull
same-product views together and push different products apart will
produce a retrieval signal that is qualitatively different from, and
stronger than, one derived from a general-purpose pre-trained model.
All metrics in Tables~\ref{tab:recall}--\ref{tab:map} are averaged over
three independent seeds corresponding to team roll numbers 32, 75, and 78.

\paragraph{Impact of contrastive fine-tuning.}
The most immediate observation is that fine-tuned visual embeddings alone
($\alpha=1$) decisively outperform Condition~A on every metric at every $K$.
Recall@15 rises from $0.6093$ to $0.6578$, NDCG@15 from $0.6428$ to
$0.7446$, and mAP@15 from $0.3394$ to $0.7969$, an absolute gain of
$+0.4575$ on the ranking quality metric.

The asymmetry between Recall and mAP gains is informative.
Recall@K is binary: it measures only whether at least one correct item
appears in the top-$K$, regardless of where.
Recall improving by $+0.049$ while mAP improves by $+0.457$ tells us
that fine-tuning does not primarily broaden coverage; it 
rearranges the ranking so that the correct item moves from a mid-list
position to near rank~1.
This is the mechanism InfoNCE is built for: with a batch of 256 images
drawn from diverse \texttt{item\_id}s, every pair of distinct products
serves as a negative example, so the encoder is continuously penalised
for placing the wrong item at rank~1.
Over training this forces the vision encoder to resolve the fine-grained
visual distinctions (collar geometry, hem cut, stitch density) that the
pre-trained CLIP model, optimised for broad image-text alignment,
does not need to resolve.

\paragraph{Why equal-weight text fusion ($\alpha=0.5$) degrades performance.}
Introducing text embeddings at equal weight relative to the visual
embeddings makes every metric worse than vision-only.
NDCG@5 falls from $0.6185$ to $0.5856$; mAP@15 falls from $0.7969$ to
$0.7364$.
The explanation is a representational mismatch between the two encoders.
The fine-tuned vision encoder has been sculpted to be
\emph{instance-discriminative}: two images of the same product cluster
tightly regardless of pose or background variation.
The BLIP-2 text encoder, which remains frozen throughout training, was
never optimised for this task.
Its output for ``a white linen shirt with button collar'' is essentially
identical for every white linen shirt in the gallery, irrespective of
which specific product it belongs to.
Fusing 50\% of this category-level signal into instance-discriminative
vectors smears the precise boundaries the visual encoder has learned,
reducing the separability of products within the same semantic category.

\paragraph{Why image-dominant fusion ($\alpha=0.7$) improves over both alternatives.}
At $\alpha=0.7$ every single metric improves over both $\alpha=1$ and
$\alpha=0.5$: Recall@15 reaches $0.6878$, NDCG@15 reaches $0.7853$,
and mAP@15 reaches $0.8411$.
The gain over pure vision-only ($\alpha=1$) shows that text does carry
complementary information, but only when it is a minority contributor.
At 30\% weight, the caption embedding provides a coarse semantic anchor
that corrects cases where two visually close products belong to different
broad categories, without overriding the fine-grained instance geometry.
The text acts as a regulariser on semantic coherence, not a co-equal
retrieval signal.
This interaction only arises because the two sources have been trained
for fundamentally different objectives: the visual encoder for instance
identity, the text encoder for semantic category membership.
The optimal $\alpha$ is therefore not a hyperparameter to tune blindly;
it reflects the relative granularity at which each modality operates.

\paragraph{Seed stability.}
The $\alpha=0.7$ results are also the most consistent across seeds
(mAP@15 std $= 0.0057$).
The $\alpha=0.5$ configuration shows nearly three times the variance
(std $= 0.0167$) on the same metric.
Higher variance at equal weight is expected: when the text signal
contributes as much as the visual signal, the outcome depends more
sensitively on the exact checkpoint the visual encoder converges to
under each seed, since the two components do not reinforce each other.

% ─────────────────────────────────────────────────────────────────────────────
\section{Interactive Demo Application}
% ─────────────────────────────────────────────────────────────────────────────

We develop a Streamlit-based interactive web application demonstrating the
complete retrieval pipeline.

The workflow is as follows:

\begin{enumerate}[noitemsep]
    \item User uploads a clothing image.
    \item YOLO detects and crops the product region.
    \item Original and cropped images are displayed.
    \item User confirms the crop.
    \item CLIP generates the query embedding.
    \item HNSW retrieves top-$K$ matching products.
    \item Retrieved products are displayed with similarity scores.
\end{enumerate}

The gallery embeddings are precomputed offline and cached during application
startup for fast inference. Additionally, a batch evaluation script is provided
in the repository to compute project metrics end-to-end on a folder of query images.

% ─────────────────────────────────────────────────────────────────────────────
\section{Conclusion}
% ─────────────────────────────────────────────────────────────────────────────

This project presents a complete multimodal visual product retrieval system
for fashion search.

The experiments demonstrate that:

\begin{itemize}[noitemsep]
    \item YOLO improves localization by removing background clutter.
    \item BLIP-2 semantic captions improve multimodal alignment.
    \item CLIP contrastive fine-tuning significantly improves retrieval
          quality.
    \item HNSW enables scalable and low-latency nearest-neighbor retrieval.
\end{itemize}

The ablation study confirms that combining semantic fusion with task-specific
fine-tuning produces the best retrieval performance.

% Future work could explore:

% \begin{itemize}[noitemsep]
%     \item stronger caption generation,
%     \item full encoder fine-tuning,
%     \item larger retrieval indexes,
%     \item and transformer-based re-ranking methods.
% \end{itemize}

% ─────────────────────────────────────────────────────────────────────────────
\section*{GitHub Repository}
% ─────────────────────────────────────────────────────────────────────────────

All code, models, and scripts are available at:

\url{https://github.com/Arnav-Oruganty/VR-Final-Project}

\end{document}